
\documentclass[12pt]{article}
%\usepackage[libertine]{newtxmath}
\usepackage{geometry}
\usepackage{caption}
\usepackage{titling}
\usepackage{xcolor}
\usepackage{amsmath,amscd}
\usepackage{stmaryrd}
\usepackage{caption}
\usepackage{hyperref}
\usepackage{graphicx}
\usepackage{algorithm}
\usepackage{algpseudocode}
\usepackage{fontspec}
\usepackage{unicode-math}

\input{macros.tex}

\geometry{a4paper, top=2.5cm, bottom=2.5cm, left=2.2cm, right=2.2cm}
\setlength{\droptitle}{-3cm}
\hypersetup{colorlinks=true, linkcolor=blue, urlcolor=blue, citecolor=blue}
\captionsetup{
    justification=centering,
    font=small,
}

%\setmainfont{TeX Gyre Pagella}
%\setmathfont{Asana Math}

\setmainfont{Libertinus Serif}
\setsansfont{Libertinus Sans}
%\setmathfont{Libertinus Math}
\setmathfont{STIX Two Math}
%\setmonofont{Libertinus Mono}
\setmonofont{lmmono10-regular.otf}[
  BoldFont = lmmonolt10-bold.otf,
  ItalicFont = lmmono10-italic.otf,
]

\title{Mathematical Definitions for ML-KEM}
\AuthorDeviPrasad
\date{\today}

\begin{document}
\maketitle

\newpage
\begin{tabular}[H]{r p{0.85\linewidth}}
$n$            & An integer value 256.\SixPtSp

$q$            & Denotes a prime integer $3329 = 2^8 \cdot 13 + 1$.\SixPtSp

$\zeta$        & Denotes the integer value 17.\\
               & This is the n-th root of unity modulo $q$ where $n = 256$.\\
               & In other words, $\zeta^{256} = 1 \Mod{q}$. \TwoPtSp
               & \SmallBlackTrRt Python Snippet 1:\OnePtSp
               & \quad\quad\texttt{{assert pow(17, 256, 3329) == 1}}\SixPtSp

$\mathbb{Z}$   & The set of integers.\SixPtSp

$\mathbb{R}$   & The set of real numbers.\SixPtSp

$\mathbb{Z}_q$ & The \textbf{\textit{field}} of integers modulo $q$.\\
               & The underlying set is \{0, 1, \ldots, 3328\}\\
               & This set admits addition and multiplication operations modulo $q$.\SixPtSp

$\mathbb{Z}_m$  & The (quotient) \textbf{\textit{ring}} of integers modulo arbitrary $m$.\\
                & The underlying set is \{0, 1, \ldots, m-1\}\\
                & This set admits addition and multiplication operations modulo $m$.\SixPtSp

$\mathbb{Z}^{n}_{m}$
               & The set of $n$-tuples over $\mathbb{Z}_m$, equipped with $\mathbb{Z}_m$-module structure.\\
               & It is a set of vectors/tuples of length $n$.\\
               & Each element of the vector/tuple is in $Z_m$.\TwoPtSp
               & \SmallBlackTrRt Example 1:\OnePtSp
               & \quad\quad $\{(0,3328,1),(1665,0,802),(3328,3328,3328),(100,200,300)\} \in \mathbb{Z}^{3}_{q}$ \TwoPtSp
               & \SmallBlackTrRt Counterexample 1:\OnePtSp
               & \quad\quad $\{(0,\Red{3329}),(\Red{65535},\Red{8802}),(0,0),(29,\Red{3330}),(100,200)\} \not{\in} \mathbb{Z}^{2}_{q}$
               \SixPtSp


$\mathbb{Z}_{q}[X]$
               & The ring of polynomials of an arbitrary degree with coefficients in $\mathbb{Z}_q$.\\
               & Addition and multiplication are usual polynomial operations.\\
               & The coefficients are reduced modulo $q$.\\
               & It is a polynomial ring over a field - coefficients are members of $\mathbb{Z}_q$.\SixPtSp


$f \in \mathbb{Z}_q$
               & A polynomial of the form $f = a_0 + a_1 X + a_2 X + \ldots + a_{m} X^{m}$,
                 where $a_j \in \mathbb{Z}_q$, and $m$ is arbitrary.\SixPtSp

$R_q$          & The ring $Z_q[X]/(X^n+1)$ consists of polynomials of the form $f$.\\
               & This ring admits addition and multiplication of polynomials modulo $X^n + 1$. \\
               & $f \in R_q$ is a polynomial in the ring $R_q$.\SixPtSp

$f \in R_q$    & is a polynomial of the form $f = a_0 + a_1 X + a_2 X + \ldots + a_{255} X^{255}$
                 where $a_j \in \mathbb{Z}_q$.\SixPtSp
\end{tabular}


In FIPS 203, $q$ is prime, and therefore $Z_q$ is a \textit{field}, i.e., every element has a multiplicative inverse.

\newpage

\section{Algebraic Structures}

\subsection{Additive Abelian Group}

\subsection{Rings}
A ring $(R, +, \times)$ consists of a set R with two binary operations $+$ and $\times$ on R. These two
operations satisfy the following axioms:

\begin{enumerate}
    \item $(R, +)$ is an abelian group with identity denoted 0.
    \item There is a multiplicative identity, denoted by 1, such that $1 \times a = a \times 1$
              for all $a \in R$.
    \item Additive identity 1 is different from multiplicative identity 0
        \begin{itemize}
            \item $1 \neq 0$.
        \end{itemize}
    \item The operation  $\times$ is associative
        \begin{itemize}
            \item $a \times (b \times c) = (a \times b) \times c$ for all $a, b, c \in R.$
        \end{itemize}
    \item The operation $\times$ is distributive over +
        \begin{itemize}
            \item $a \times (b + c) = (a \times b) + (a \times c)$ for all $a, b, c \in R.$
        \end{itemize}
\end{enumerate}

\subsection{Polynomial Rings}

\subsection{Vector Space and Vector Basis}
\noindent
A \textbf{vector space} $V$ is a set that is closed under vector addition and scalar multiplication.

\noindent
In general, the scalars are members of a \textit{field} $F$, and we say something like "vector space over a filed".

\noindent
Let $V$ be a vector space over a field $F$ (for example, over $\mathbb{R}$).


\noindent
Let $v_1, v_2, \ldots, v_n \in V$ be linearly independent vectors in $V$. Then, the \textbf{span}
of these vectors is the set of all linear combinations of the vectors.
$$
    \textit{span}_{\mathbb{R}}(v_1, v_2, \ldots, v_n) = \{a_1 v_1 + a_2 v_2 + \ldots + a_n v_n \mid a_i \in \mathbb{R}\}
$$

\noindent
The \textbf{Z-linear span} of these vectors is the set of all linear combinations
of these vectors with integer coefficients.
$$
    \textit{span}_{\mathbb{Z}}(v_1, v_2, \ldots, v_n) = \{a_1 v_1 + a_2 v_2 + \ldots + a_n v_n \mid a_i \in \mathbb{Z}\}
$$

\noindent
\subsection{Module}
A \textbf{module} is similar to vector space, but its coefficients are members of a \text{ring}, instead of a \text{field}. Therefore, we say "modules generalize the concept of a vector space".

\subsection{Lattice}


\subsection{Module Lattice}

\[
\begin{array}{ll}
    \text{LWE} & \text{vectors over} \ \mathbb{Z}_q \TwoPtSp
    \text{Ring-LWE} & \text{module rank 1 over} \ R_q \TwoPtSp
    \text{Module-LWE} & \text{module rank} \ $k$ \ \text{over} \ R_q \TwoPtSp
\end{array}
\]

\end{document}